\documentclass[11pt,colorlinks]{charlie}
\usepackage{booktabs,array}
\newcolumntype{L}[1]{>{\raggedright\arraybackslash}p{#1}}
\setlength{\emergencystretch}{3em}
\usepackage[capitalize,nameinlink]{cleveref}
\usepackage{aliascnt}
\makeatletter
\def\ch@redefthm#1#2#3{\expandafter\let\csname #1\endcsname\relax
  \expandafter\let\csname end#1\endcsname\relax
  \newaliascnt{#1}{theorem}\newtheorem{#1}[#1]{#2}\aliascntresetthe{#1}%
  \crefname{#1}{#2}{#3}\Crefname{#1}{#2}{#3}}
\crefname{theorem}{Theorem}{Theorems}\Crefname{theorem}{Theorem}{Theorems}
\theoremstyle{chplain}
\ch@redefthm{lemma}{Lemma}{Lemmas}
\ch@redefthm{proposition}{Proposition}{Propositions}
\ch@redefthm{corollary}{Corollary}{Corollaries}
\ch@redefthm{conjecture}{Conjecture}{Conjectures}
\theoremstyle{chdefinition}
\ch@redefthm{definition}{Definition}{Definitions}
\ch@redefthm{example}{Example}{Examples}
\ch@redefthm{hypothesis}{Hypothesis}{Hypotheses}
\makeatother

\newcommand{\R}{\mathbb{R}}
\newcommand{\C}{\mathbb{C}}
\newcommand{\Z}{\mathbb{Z}}
\newcommand{\Q}{\mathbb{Q}}
\newcommand{\N}{\mathbb{N}}
\newcommand{\E}{\mathbb{E}}
\newcommand{\one}{\mathbf{1}}
\newcommand{\dd}{\,\mathrm{d}}

\newif\ifdraftmarkers \draftmarkersfalse
\newcounter{gapcount}
\newcommand{\gap}[1]{\ifdraftmarkers\stepcounter{gapcount}%
  \par\smallskip\noindent\fcolorbox{red!60!black}{red!5}{\parbox{\dimexpr\linewidth-2\fboxsep-2\fboxrule}{%
  \textcolor{red!60!black}{\textbf{Gap \thegapcount.}} #1}}\par\smallskip\fi}
\newcommand{\unchecked}[1]{\ifdraftmarkers%
  \par\smallskip\noindent\fcolorbox{orange!80!black}{orange!6}{\parbox{\dimexpr\linewidth-2\fboxsep-2\fboxrule}{%
  \textcolor{orange!80!black}{\textbf{Unchecked.}} #1}}\par\smallskip\fi}
\newcommand{\source}[1]{\ifdraftmarkers{\textcolor{gray}{\footnotesize[source: \texttt{#1}]}}\fi}
\newcommand{\todo}[1]{\ifdraftmarkers{\textcolor{blue!70!black}{\footnotesize[\textbf{todo:} #1]}}\fi}

\newcommand{\coderepo}{\url{https://github.com/charlieoneill11/maths-releases}}
\newcommand{\aiusestatement}{The mathematical arguments and computer programs in this paper were developed by Astra (GPT-6.2 Astra, OpenAI). Fable (Fable~5.2, Anthropic) worked throughout as a second model: it proposed alternative approaches, reviewed arguments and raised objections. The manuscript was prepared from the project's research records with Claude (Anthropic). All computations were re-run before release, and the programs and certificates are available at \coderepo.}
\newcommand{\responsibilitystatement}{Charles O'Neill has checked the arguments and takes full responsibility for the correctness of the results.}
\newcommand{\mscandkeywords}[2]{\par\medskip\noindent{\small\textit{2020 Mathematics Subject Classification.} #1\par\noindent\textit{Keywords.} #2\par}}

\usepackage{amssymb}

\DeclareMathOperator{\Aut}{Aut}
\DeclareMathOperator{\Inn}{Inn}
\DeclareMathOperator{\Out}{Out}
\DeclareMathOperator{\Soc}{Soc}
\DeclareMathOperator{\Sym}{Sym}
\newcommand{\gen}[1]{\langle #1\rangle}
\renewcommand{\aiusestatement}{The argument in this paper was developed by Astra (GPT-6.2 Astra, OpenAI) from Charles O'Neill's instructions. The manuscript was prepared from the project's research records with Claude (Anthropic). The paper contains no computations.}

\title{Conjugate pairs with the full prime spectrum in finite groups}
\author{Charles O'Neill, Max Kirkby and Astra}
\date{25 November 2026}

\begin{document}
\maketitle

\begin{abstract}
For a finite group $X$ let $\pi(X)$ be the set of primes dividing $|X|$. We prove that every finite group $G$ contains two conjugate elements $a$ and $b$ with $\pi(\langle a,b\rangle)=\pi(G)$. This answers Question~17.125 of the Kourovka Notebook, a conjecture of Shumyatsky. Equivalently, every finite group in which no proper subgroup has the same prime spectrum is generated by two conjugate elements. Maslova and Revin proved this for groups whose maximal subgroups are all Hall subgroups. We study a minimal counterexample. A lifting argument through an abelian minimal normal subgroup shows that its soluble radical is trivial. For each minimal normal subgroup $S^k$ we then lift a conjugate generating pair of the quotient by the socle so that the subgroup it generates has order divisible by $|S|$. When $k=1$ this uses the theorem of Burness, Guralnick and Harper that a finite group whose proper quotients are all cyclic is generated by two conjugate elements. When $k\ge2$ it uses a choice of automorphism labels along a spanning tree of the permutation action, together with the theorems of Gasch\"utz and of Lucchini and Menegazzo. The proof depends on the classification of finite simple groups. We have also checked in Lean~4 that the theorem follows from four of the cited published results.
\end{abstract}
\mscandkeywords{20D60; 20F05, 20D06, 20E22}{prime spectrum, conjugate elements, generation of finite groups, Kourovka Notebook, minimal normal subgroups, wreath products}

\section*{AI use and author contributions}
\aiusestatement

Charles O'Neill proposed the problem and asked for it to be answered. \responsibilitystatement

\section{Introduction}
For a finite group $X$ write $\pi(X)$ for the set of prime divisors of $|X|$, the \emph{prime spectrum} of $X$. Shumyatsky conjectured that every finite group $G$ contains a pair of conjugate elements $a,b$ with $\pi(\gen{a,b})=\pi(G)$. The conjecture appears as Question~17.125 in the Kourovka Notebook \cite{Kourovka}; see \cite{MR13}.

Following Maslova and Revin, call a finite group $P$ \emph{prime-spectrum minimal} if $\pi(U)\ne\pi(P)$ for every proper subgroup $U<P$. Every finite group contains a prime-spectrum-minimal subgroup with the same prime spectrum, namely a subgroup of least order with that spectrum, and a pair of elements of a prime-spectrum-minimal group whose subgroup has the full spectrum must generate the group (\Cref{lem:reduction}). So Shumyatsky's conjecture is equivalent to the statement that every prime-spectrum-minimal group is generated by two conjugate elements. Maslova and Revin \cite{MR13} proved this for groups all of whose maximal subgroups are Hall subgroups, and Maslova \cite{Maslova} proved it for prime-spectrum-minimal groups whose nonabelian composition factors all lie in $\{\mathrm{PSL}_2(7),\mathrm{PSL}_2(11),\mathrm{PSL}_5(2)\}$. We prove the conjecture in general.

\begin{theorem}\label{thm:main}
Every finite group $G$ contains two conjugate elements $a$ and $b$ with $\pi(\gen{a,b})=\pi(G)$.
\end{theorem}

\begin{corollary}\label{cor:psm}
Every prime-spectrum-minimal finite group is generated by two conjugate elements.
\end{corollary}

\paragraph{Outline.}
By \Cref{lem:reduction} it suffices to prove \Cref{cor:psm}, and we consider a counterexample $G$ of least order. Prime-spectrum minimality passes to quotients, so every proper quotient of $G$ is generated by two conjugate elements. \Cref{sec:radical} shows that the soluble radical of $G$ is trivial: lifting a conjugate generating pair through an abelian minimal normal subgroup $N$ either generates $G$ or produces a complement to $N$, and in the second case the pair can be modified so that it generates $G$. Then the socle $M$ of $G$ is a direct product of minimal normal subgroups $N_i\cong S_i^{k_i}$ with $S_i$ nonabelian simple. We fix a conjugate generating pair of $G/M$ and choose its lifts block by block. Since $G$ is prime-spectrum minimal, it is enough that the subgroup generated by the lifts has order divisible by $|G/M|$ and by every $|S_i|$; we never need to show directly that it is all of $G$.

For a block with $k_i=1$ (\Cref{sec:almostsimple}) the image of $G$ in $\Aut(S_i)$ is an almost simple prime-spectrum-minimal group. Its outer part has order prime to $|S_i|$ and is therefore cyclic, so every proper quotient is cyclic and the theorem of Burness, Guralnick and Harper \cite{BGH} gives a conjugate generating pair, which we adjust to have the prescribed images. For a block with $k_i\ge2$ (\Cref{sec:wreath}) we work in $\Aut(S^k)=\Aut(S)\wr\Sym(k)$. The equation $y=gxg^{-1}$ can be solved for $g$ after the labels of $x$ and of all but one edge on each cycle of $y$ have been chosen freely. Two of the free labels lie off a spanning tree of the permutation action, so the automorphisms of the first factor induced by two closed walks can be prescribed independently within their cosets of $\Inn(S)$. The lifting lemma of Gasch\"utz \cite{Gaschutz}, together with the theorem of Lucchini and Menegazzo \cite{LM}, lets us choose these two automorphisms to generate a group containing $\Inn(S)$. \Cref{sec:proof} assembles the blocks.

The proof uses the classification of finite simple groups through Schreier's conjecture, the theorem on coprime automorphisms of simple groups \cite[Theorem~5.1]{SV} and the results of \cite{BGH}. It contains no computations. \Cref{sec:lean} describes a Lean~4 proof that \Cref{thm:main} follows from four of the published results we cite.

\paragraph{Notation.}
All groups are finite. For elements $g,x$ of a group we write $gxg^{-1}$ for the conjugate of $x$ by $g$, and $\gen{X}$ for the subgroup generated by $X$. We write $|x|$ for the order of $x$, $C_X(Y)$ for a centralizer, $d(X)$ for the least number of generators of $X$, and $R(X)$ for the soluble radical of $X$, its largest soluble normal subgroup.

\section{Preliminaries}\label{sec:prelim}

\subsection{Prime-spectrum-minimal groups}\label{sec:reduction}

\begin{lemma}\label{lem:reduction}
\begin{enumerate}
\item Let $F$ be a finite group and let $P\le F$ be a subgroup of least order with $\pi(P)=\pi(F)$. Then $P$ is prime-spectrum minimal.
\item If $P$ is prime-spectrum minimal and $a,b\in P$ satisfy $\pi(\gen{a,b})=\pi(P)$, then $\gen{a,b}=P$.
\item If $P$ is prime-spectrum minimal and $K\lhd P$, then $P/K$ is prime-spectrum minimal.
\end{enumerate}
In particular, \Cref{thm:main} is equivalent to \Cref{cor:psm}.
\end{lemma}

\begin{proof}
(a) A proper subgroup $U<P$ with $\pi(U)=\pi(P)=\pi(F)$ would have smaller order than $P$.

(b) This is the definition, applied to $U=\gen{a,b}$.

(c) Let $K\le U\le P$ with $\pi(U/K)=\pi(P/K)$. Then
\[
 \pi(U)=\pi(K)\cup\pi(U/K)=\pi(K)\cup\pi(P/K)=\pi(P),
\]
so $U=P$ and $U/K=P/K$.

For the last assertion, suppose \Cref{cor:psm} holds and let $F$ be finite. Choose $P$ as in (a). Then $P=\gen{a,gag^{-1}}$ for some $a,g\in P$, and $a$ and $gag^{-1}$ are conjugate in $F$ with $\pi(\gen{a,gag^{-1}})=\pi(P)=\pi(F)$. Conversely, if \Cref{thm:main} holds and $P$ is prime-spectrum minimal, it gives $a,g\in P$ with $\pi(\gen{a,gag^{-1}})=\pi(P)$, and $\gen{a,gag^{-1}}=P$ by (b).
\end{proof}

\subsection{Results from the literature}

We use the following facts.

\begin{theorem}[Schur--Zassenhaus; see {\cite[Chapter~3]{Isaacs}}]\label{thm:sz}
Let $N$ be a normal subgroup of $G$ with $\gcd(|N|,|G/N|)=1$. Then $N$ has a complement in $G$. If $N$ or $G/N$ is soluble, all complements of $N$ are conjugate in $G$.
\end{theorem}

\begin{theorem}[Schreier's conjecture; see {\cite{GLS3}}]\label{thm:schreier}
If $S$ is a nonabelian finite simple group, then $\Out(S)$ is soluble.
\end{theorem}

\begin{theorem}[{\cite[Theorem~5.1]{SV}}]\label{thm:coprime}
Let $S$ be a nonabelian simple group and let $C$ be a group acting faithfully on $S$ with $\gcd(|C|,|S|)=1$. Then $C$ is cyclic.
\end{theorem}

\begin{theorem}[Gasch\"utz {\cite{Gaschutz}}]\label{thm:gaschutz}
Let $N\lhd L$ and $r\ge d(L)$. For every $r$-tuple $(\bar l_1,\dots,\bar l_r)$ generating $L/N$ there are $l_j\in\bar l_j$ such that $L=\gen{l_1,\dots,l_r}$.
\end{theorem}

\begin{theorem}[Lucchini--Menegazzo {\cite[Theorem~1.1]{LM}}]\label{thm:lm}
If $L$ is a noncyclic finite group with a unique minimal normal subgroup $N$, then $d(L)=\max\{2,d(L/N)\}$.
\end{theorem}

\begin{theorem}[Burness--Guralnick--Harper {\cite[Corollary~4(ii)]{BGH}}]\label{thm:bgh}
Let $L$ be a finite group such that every proper quotient of $L$ is cyclic. If $L$ is not isomorphic to $C_p\times C_p$ for a prime $p$, then $L$ is generated by two conjugate elements.
\end{theorem}

\subsection{Almost simple groups}
Let $S$ be a nonabelian simple group and identify $S$ with $I=\Inn(S)\lhd\Aut(S)$.

\begin{lemma}\label{lem:centralizer}
\begin{enumerate}
\item $C_{\Aut(S)}(I)=1$.
\item If $I\le L\le\Aut(S)$, then every nontrivial normal subgroup of $L$ contains $I$. In particular $I$ is the unique minimal normal subgroup of $L$.
\end{enumerate}
\end{lemma}

\begin{proof}
(a) Let $\alpha\in\Aut(S)$ commute with every inner automorphism. For $s\in S$ write $\iota_s$ for conjugation by $s$; then $\alpha\iota_s\alpha^{-1}=\iota_{\alpha(s)}$, so $\iota_{\alpha(s)}=\iota_s$ and $\alpha(s)s^{-1}$ lies in the centre of $S$, which is trivial. Thus $\alpha=1$.

(b) Let $1\ne K\lhd L$. Then $K\cap I\lhd I$, so $K\cap I$ is $1$ or $I$ because $I\cong S$ is simple. If $K\cap I=1$, then $[K,I]\le K\cap I=1$, so $K\le C_{\Aut(S)}(I)=1$ by (a), a contradiction. Hence $I\le K$.
\end{proof}

\section{The soluble radical of a minimal counterexample}\label{sec:radical}

Throughout the rest of the paper, $G$ denotes a prime-spectrum-minimal group that is not generated by two conjugate elements and has least order among such groups. By \Cref{lem:reduction}(c) and the minimality of $|G|$, every proper quotient $G/K$ with $K\ne1$ is generated by two conjugate elements. Clearly $G\ne1$.

\begin{proposition}\label{prop:radical}
$R(G)=1$.
\end{proposition}

\begin{proof}
Suppose $R(G)\ne1$. The last nontrivial term of the derived series of $R(G)$ is abelian and characteristic in $R(G)$, hence normal in $G$. Let $N$ be a minimal normal subgroup of $G$ contained in it. Then $N$ is abelian, and since its subgroup of elements of order dividing a prime $p\in\pi(N)$ is characteristic in $N$, the group $N$ is an elementary abelian $p$-group.

By the minimality of $|G|$, there are $\bar x,\bar g\in G/N$ with $G/N=\gen{\bar x,\bar g\bar x\bar g^{-1}}$. Choose $x\in\bar x$ and $g\in\bar g$, put $y=gxg^{-1}$ and $H=\gen{x,y}$. Then $HN=G$, so $|G:H|=|N:H\cap N|$ is a power of $p$ and $\pi(G)\setminus\pi(H)\subseteq\{p\}$. If $H=G$, then $G$ is generated by the conjugate elements $x,y$, a contradiction. So $H<G$, and prime-spectrum minimality gives $\pi(H)\ne\pi(G)$, that is, $p\notin\pi(H)$. Hence $H\cap N=1$, $H$ is a complement to $N$, and $p\nmid|H|=|G/N|$. In particular $N$ is a normal Hall subgroup of $G$. We obtain a contradiction in each of two cases.

\emph{Case 1: $C_N(x)\ne1$.} Choose $1\ne v\in C_N(x)$. Since $x\in H$, the order of $x$ is prime to $p$, and since $x$ and $v$ commute, $(xv)^{|x|}=v^{|x|}\ne1$. Hence $p$ divides $|xv|$. Put $H'=\gen{xv,\,g(xv)g^{-1}}$. As $v\in N$, the two generators of $H'$ map to $\bar x$ and $\bar g\bar x\bar g^{-1}$, so $H'N=G$ and $|G/N|$ divides $|H'|$. Together with $p\in\pi(H')$ this gives $\pi(H')=\pi(G)$, hence $H'=G$ by prime-spectrum minimality. So $G$ is generated by two conjugate elements, a contradiction.

\emph{Case 2: $C_N(x)=1$.} First, $H$ is the only complement to $N$ that contains $x$. Indeed, by \Cref{thm:sz} (with $N$ abelian) every complement is conjugate to $H$ in $G=NH$, hence equal to $nHn^{-1}$ for some $n\in N$. If $x\in nHn^{-1}$, then $n^{-1}xn\in H$. The elements $n^{-1}xn$ and $x$ of $H$ have the same image in $G/N$, and $H\cap N=1$, so $n^{-1}xn=x$. Thus $n\in C_N(x)=1$ and $nHn^{-1}=H$.

Next, $C_N(y)=gC_N(x)g^{-1}=1$ because $N\lhd G$. Choose $1\ne n\in N$ and put $y'=nyn^{-1}=(ng)x(ng)^{-1}$. If $y'\in H$, then $y'$ and $y$ are elements of $H$ with the same image in $G/N$, so $y'=y$ and $n\in C_N(y)=1$, which is false. Hence $y'\notin H$. Let $K=\gen{x,y'}$. Its generators map to $\bar x$ and $\bar g\bar x\bar g^{-1}$, so $KN=G$. If $K<G$, the argument of the second paragraph, applied to $K$ in place of $H$, shows that $K$ is a complement to $N$. It contains $x$, so $K=H$ by the uniqueness just proved, which contradicts $y'\notin H$. Therefore $K=G$, and $G$ is generated by the conjugate elements $x$ and $y'$, a contradiction.
\end{proof}

\Cref{prop:radical} is in the spirit of the structural reductions of Maslova and Revin \cite{MR13}.

\section{The socle and its permutation action}\label{sec:socle}

Let $M=\Soc(G)$ be the product of the minimal normal subgroups of $G$. Since $R(G)=1$ by \Cref{prop:radical}, every minimal normal subgroup of $G$ is nonabelian. A nonabelian minimal normal subgroup is characteristically simple, so it is isomorphic to $S^k$ for a nonabelian simple group $S$ and some $k\ge1$; its $k$ simple direct factors are exactly its minimal normal subgroups, and they are permuted by conjugation.

\begin{lemma}\label{lem:socle}
Let $N_1,\dots,N_r$ be the minimal normal subgroups of $G$, with $N_i\cong S_i^{k_i}$.
\begin{enumerate}
\item $M=N_1\times\dots\times N_r$, and $|G|=|G/M|\prod_i|S_i|^{k_i}$.
\item For each $i$, $G$ acts transitively by conjugation on the $k_i$ simple factors of $N_i$, and $M$ acts trivially on them.
\end{enumerate}
\end{lemma}

\begin{proof}
(a) For $i\ne j$ we have $[N_i,N_j]\le N_i\cap N_j=1$, so distinct minimal normal subgroups commute. Each $N_i$ has trivial centre. If $N_i\cap\prod_{j\ne i}N_j\ne1$, then this intersection, being normal in $G$, contains $N_i$ by minimality, so $N_i$ is centralized by itself and is abelian, a contradiction. Hence the product is direct.

(b) The product of the simple factors of $N_i$ in one $G$-orbit is a nontrivial normal subgroup of $G$ contained in $N_i$, so by minimality it is $N_i$ and the orbit contains all the factors. An element of $N_i$ normalizes each simple factor of $N_i$, since the factors are normal in $N_i$, and $N_j$ centralizes $N_i$ for $j\ne i$. So $M$ fixes every factor.
\end{proof}

Since $M\ne1$, the quotient $G/M$ is prime-spectrum minimal of smaller order, so there are $x_0,g_0\in G$ with
\begin{equation}\label{eq:quotient-pair}
 G/M=\gen{x_0M,\ y_0M},\qquad y_0=g_0x_0g_0^{-1}.
\end{equation}
We fix them for the rest of the paper. For each $i$ let $\rho_i\colon G\to\Aut(N_i)$ be the conjugation action. Then $\rho_i(N_i)=\Inn(N_i)$ and $\rho_i(N_j)=1$ for $j\ne i$, so $\rho_i(M)=\Inn(N_i)$.

\begin{lemma}\label{lem:transitive}
Let $x\in x_0M$ and $g\in g_0M$, and put $y=gxg^{-1}$. Then $y\in y_0M$, and for each $i$ the permutations induced by $x$ and $y$ on the simple factors of $N_i$ generate a transitive group.
\end{lemma}

\begin{proof}
As $M\lhd G$, $gxg^{-1}\in g_0x_0g_0^{-1}M$. By \Cref{lem:socle}(b) the permutation action of $G$ on the factors of $N_i$ factors through $G/M$, which is generated by the images of $x$ and $y$, and it is transitive.
\end{proof}

\section{Blocks with one simple factor}\label{sec:almostsimple}

\begin{proposition}\label{prop:almostsimple}
Let $S$ be a nonabelian simple group, identified with $\Inn(S)$, and let $S\le L\le\Aut(S)$ with $L$ prime-spectrum minimal. Then $L/S$ is cyclic. Moreover, if $c,d\in L/S$ and $c$ generates $L/S$, then there are $a\in c$ and $h\in d$ with $L=\gen{a,hah^{-1}}$.
\end{proposition}

\begin{proof}
By \Cref{thm:schreier}, $L/S\le\Out(S)$ is soluble. By Hall's theorem (see \cite[Chapter~3]{Isaacs}), $L/S$ has a Hall subgroup $K/S$ for the set of primes outside $\pi(S)$. Then
\[
 \pi(K)=\pi(S)\cup\pi(K/S)\supseteq\pi(S)\cup\bigl(\pi(L/S)\setminus\pi(S)\bigr)=\pi(L),
\]
so $K=L$ by prime-spectrum minimality, and $\gcd(|L/S|,|S|)=1$. By \Cref{thm:sz}, $S$ has a complement $C$ in $L$. As a subgroup of $\Aut(S)$, the group $C$ acts faithfully on $S$, and its order is prime to $|S|$. By \Cref{thm:coprime}, $C$ is cyclic, and so is $L/S\cong C$.

By \Cref{lem:centralizer}(b), every nontrivial normal subgroup of $L$ contains $S$, so every proper quotient of $L$ is a quotient of $L/S$ and is cyclic. Since $L$ contains the nonabelian group $S$, it is not isomorphic to $C_p\times C_p$. By \Cref{thm:bgh} there are $a_1,g_1\in L$ with $L=\gen{a_1,b_1}$, where $b_1=g_1a_1g_1^{-1}$.

It remains to adjust the cosets. Since $L/S$ is abelian, $b_1S=a_1S$, so $a_1S$ generates $L/S$. Put $q=|L/S|$ and $n=|a_1|$. As $c$ also generates the cyclic group $L/S$ of order $q$, we have $c=(a_1S)^r$ for some integer $r$ prime to $q$. By the Chinese remainder theorem choose an integer $e$ with
\[
 e\equiv r\pmod q,\qquad e\equiv1\pmod\ell\ \text{ for every prime }\ell\mid n\text{ with }\ell\nmid q .
\]
Every prime divisor of $n$ either divides $q$, and then does not divide $e$ because $\gcd(r,q)=1$, or does not divide $q$, and then does not divide $e$ by construction. Hence $\gcd(e,n)=1$. (When $q=1$ the first congruence is empty.) Put $a=a_1^e$. Then $\gen{a}=\gen{a_1}$ and $\gen{b_1^e}=\gen{b_1}$, because $|b_1|=n$ as well, so $L=\gen{a,b_1^e}$. Also $aS=(a_1S)^e=(a_1S)^r=c$ and $b_1^e=g_1ag_1^{-1}$.

Finally choose an integer $j$ with $c^j=(g_1S)^{-1}d$, which is possible because $c$ generates $L/S$, and put $h=g_1a^j$. Then $hS=(g_1S)c^j=d$ and $hah^{-1}=g_1ag_1^{-1}=b_1^e$. Hence $L=\gen{a,hah^{-1}}$ with $a\in c$ and $h\in d$.
\end{proof}

\section{Blocks with several simple factors}\label{sec:wreath}

In this section $S$ is a nonabelian simple group, $k\ge2$, $I=\Inn(S)$, and $W=\Aut(S^k)$. We write $B=\Inn(S^k)\lhd W$.

\subsection{Coordinates}
The minimal normal subgroups of $S^k$ are its $k$ simple factors, so every automorphism of $S^k$ permutes them. Number the factors $1,\dots,k$ and identify each with $S$. Then an element $w\in W$ is given by a permutation $\hat w$ of $\{1,\dots,k\}$ and labels $w_1,\dots,w_k\in\Aut(S)$: it maps the element $s$ of factor $i$ to the element $w_i(s)$ of factor $\hat w(i)$. Conversely every such pair defines an automorphism. If $u,v\in W$ and $v$ is applied first, then
\begin{equation}\label{eq:compose}
 \widehat{uv}=\hat u\hat v,\qquad (uv)_i=u_{\hat v(i)}\,v_i\qquad(1\le i\le k),
\end{equation}
and consequently $(w^{-1})_{\hat w(i)}=w_i^{-1}$. The subgroup $B$ consists of the elements with trivial permutation and all labels in $I$. By \eqref{eq:compose}, for $w\in W$ the coset $wB$ consists of all elements with permutation $\hat w$ and labels $z_i\in w_iI$, chosen independently for each $i$. The stabilizer $W_1$ of factor $1$ maps to $\Aut(S)$ by $w\mapsto w_1$, and this is a homomorphism by \eqref{eq:compose}.

For $w\in W$ and $m\ge1$ we write $(w^m)_i=w_{\hat w^{m-1}(i)}\cdots w_{\hat w(i)}w_i$, which agrees with \eqref{eq:compose}. All statements below about images in $\Out(S)=\Aut(S)/I$ use that $I$ is normal in $\Aut(S)$.

\subsection{Solving the conjugacy equation}

Fix $\xi,\gamma\in W$ and put $\eta=\gamma\xi\gamma^{-1}$. Write $\sigma=\hat\xi$, $\tau=\hat\eta$ and $\beta=\hat\gamma$, so that $\tau=\beta\sigma\beta^{-1}$ and $\beta$ maps each $\sigma$-cycle onto a $\tau$-cycle of the same length, distinct $\sigma$-cycles to distinct $\tau$-cycles.

\begin{lemma}\label{lem:closure}
In each $\tau$-cycle $Z$ choose a position $r_Z\in Z$, called reserved. Let $x\in\xi B$, and prescribe labels $y_j\in\eta_jI$ for every $j$ that is not reserved. Then there are labels $y_{r_Z}\in\eta_{r_Z}I$ at the reserved positions and an element $g\in\gamma B$ such that the element $y$ with permutation $\tau$ and labels $(y_j)$ satisfies $y=gxg^{-1}$. In particular $y\in\eta B$.
\end{lemma}

\begin{proof}
We look for $g$ with permutation $\beta$. Since $\beta\sigma=\tau\beta$, the equation $gx=yg$ holds for the permutations, and by \eqref{eq:compose} it holds for the labels if and only if
\begin{equation}\label{eq:recurrence}
 g_{\sigma(i)}=y_{\beta(i)}\,g_i\,x_i^{-1}\qquad(1\le i\le k).
\end{equation}
Let $O$ be a $\sigma$-cycle of length $m$, choose a base point $i\in O$, and put $g_i=\gamma_i$. Using \eqref{eq:recurrence} to define $g$ successively at $\sigma(i),\dots,\sigma^{m-1}(i)$, induction and $\beta\sigma^t=\tau^t\beta$ give
\[
 g_{\sigma^t(i)}=(y^t)_{\beta(i)}\,g_i\,(x^t)_i^{-1}\qquad(0\le t\le m),
\]
where $y^t$ is computed from the labels $y_j$ and the permutation $\tau$. The relation \eqref{eq:recurrence} at the last point $\sigma^{m-1}(i)$ of the cycle is consistent if and only if the value for $t=m$ equals $g_i$, that is,
\begin{equation}\label{eq:closure}
 (y^m)_{\beta(i)}=g_i\,(x^m)_i\,g_i^{-1}.
\end{equation}
The left-hand side is the ordered product of the $m$ labels on the $\tau$-cycle $Z=\beta(O)$, starting at $\beta(i)$. The reserved label $y_{r_Z}$ occurs in it exactly once, so \eqref{eq:closure} has the form $P\,y_{r_Z}\,Q=R$, where $P$ and $Q$ are products of prescribed labels and $R$ is known. Set $y_{r_Z}=P^{-1}RQ^{-1}$. Doing this for every $\sigma$-cycle defines every reserved label exactly once, since each $\tau$-cycle is $\beta(O)$ for exactly one $O$. With these labels, \eqref{eq:recurrence} defines $g$ consistently on every $\sigma$-cycle, and then $gx=yg$.

It remains to check the cosets. The triple $(\xi,\gamma,\eta)$ satisfies $\gamma\xi=\eta\gamma$, so it satisfies \eqref{eq:recurrence} and \eqref{eq:closure} with $x,g,y$ replaced by $\xi,\gamma,\eta$. Modulo $I$, the prescribed labels of $y$, the labels of $x$ and the base values $g_i=\gamma_i$ agree with those of $\eta$, $\xi$ and $\gamma$. Hence modulo $I$ the reserved label $y_{r_Z}$ and $\eta_{r_Z}$ solve the same equation $\bar P\,\bar z\,\bar Q=\bar R$ in $\Out(S)$, whose solution is unique, so $y_{r_Z}\in\eta_{r_Z}I$. Then \eqref{eq:recurrence} modulo $I$ defines the same values for $g$ and for $\gamma$ along each cycle, so $g_j\in\gamma_jI$ for every $j$, that is, $g\in\gamma B$. Finally $y=gxg^{-1}\in\gamma\xi\gamma^{-1}B=\eta B$.
\end{proof}

\subsection{A spanning tree and two free labels}

Let $\Gamma$ be the directed multigraph on $\{1,\dots,k\}$ with an \emph{$x$-edge} from $i$ to $\sigma(i)$ and a \emph{$y$-edge} from $j$ to $\tau(j)$ for every $i$ and $j$. Loops and parallel edges are kept. Trees and paths below are taken in the underlying undirected multigraph. Each $\tau$-cycle of length $\ell$ gives a cycle in this multigraph: a loop if $\ell=1$, two parallel edges if $\ell=2$, and a cycle of length $\ell$ if $\ell\ge3$. The same holds for $\sigma$-cycles.

\begin{lemma}\label{lem:tree}
Suppose that $\gen{\sigma,\tau}$ is transitive on $\{1,\dots,k\}$. Then there are a spanning tree $T$ of $\Gamma$, a set $\mathcal R$ of $y$-edges containing exactly one edge of each $\tau$-cycle and no edge of $T$, and two distinct edges $e_1,e_2$ of $\Gamma$ that lie neither in $T$ nor in $\mathcal R$.
\end{lemma}

\begin{proof}
Transitivity means that $\Gamma$ is connected. Let $c$ be the number of $\sigma$-cycles. Delete one $x$-edge from each $\sigma$-cycle (the loop, for a fixed point). The remaining $x$-edges form a forest $F$ whose components are the $c$ cycles, each with an edge removed. Each deleted $x$-edge joins two vertices of one component of $F$, so a spanning tree $T\supseteq F$ of the connected multigraph $\Gamma$ consists of $F$ together with $c-1$ $y$-edges. Since a tree contains no cycle, every $\tau$-cycle has an edge outside $T$; let $\mathcal R$ contain one such edge from each $\tau$-cycle.

If $c\ge2$, let $e_1,e_2$ be two of the deleted $x$-edges. If $c=1$, then $T=F$ contains no $y$-edge, and $\tau$, being conjugate to $\sigma$, is also a single $k$-cycle. Its $k$ edges are not in $T$ and exactly one of them is in $\mathcal R$, so since $k\ge2$ at least one $y$-edge $e_2$ lies outside $T\cup\mathcal R$; let $e_1$ be the deleted $x$-edge.
\end{proof}

When $k=2$ and $c=1$, both $\sigma$ and $\tau$ are the transposition, and the two $y$-edges $1\to2$ and $2\to1$ are distinct parallel edges; the lemma then takes one as the reserved edge and the other as $e_2$.

For $x,y\in W$ with $\hat x=\sigma$ and $\hat y=\tau$, call $x_i$ the label of the $x$-edge from $i$ and $y_j$ the label of the $y$-edge from $j$. A walk in the undirected multigraph from $v_0$ to $v_n$ that traverses edges $f_1,\dots,f_n$ determines a word $\omega=\lambda_n\cdots\lambda_1$, where $\lambda_t$ is $x$ or $y$ if $f_t$ is an $x$- or $y$-edge traversed forwards, and $x^{-1}$ or $y^{-1}$ if it is traversed backwards. By \eqref{eq:compose} the element $\omega(x,y)$ maps factor $v_0$ to factor $v_n$ with label
\begin{equation}\label{eq:walk}
 \omega(x,y)_{v_0}=\mu_n\cdots\mu_1,
\end{equation}
where $\mu_t$ is the label of $f_t$ if $f_t$ is traversed forwards and its inverse otherwise. Indeed $(x^{-1})_{v}=x_{\sigma^{-1}(v)}^{-1}$, which is the inverse of the label of the $x$-edge entering $v$, and similarly for $y$.

For an edge $e$ from $s$ to $t$, let $\omega_e$ be the word of the closed walk that goes from $1$ to $s$ along $T$, traverses $e$ forwards, and returns from $t$ to $1$ along $T$. Then $\omega_e(x,y)\in W_1$, and by \eqref{eq:walk}
\begin{equation}\label{eq:walk-label}
 \omega_e(x,y)_1=A_e\,z_e\,D_e,
\end{equation}
where $z_e$ is the label of $e$ and $A_e,D_e$ are products of labels of edges of $T$ and their inverses. In particular, if $e\notin T$, then $\omega_e(x,y)_1$ depends only on the labels of $T$ and on $z_e$.

\subsection{The block lift}

\begin{proposition}\label{prop:wreath}
Let $\xi,\gamma\in W$, put $\eta=\gamma\xi\gamma^{-1}$, and suppose that $\gen{\hat\xi,\hat\eta}$ is transitive on $\{1,\dots,k\}$. Then there are $x\in\xi B$ and $g\in\gamma B$ such that, with $y=gxg^{-1}$, the order of $\gen{x,y}$ is divisible by $|S|$. More precisely, the image of $\gen{x,y}\cap W_1$ under $w\mapsto w_1$ contains $I$.
\end{proposition}

\begin{proof}
Choose $T$, $\mathcal R$, $e_1$ and $e_2$ by \Cref{lem:tree}, and take the reserved positions of \Cref{lem:closure} to be the sources of the edges in $\mathcal R$. Give every $x$-edge and every $y$-edge outside $\mathcal R\cup\{e_1,e_2\}$ the label of $\xi$ or $\eta$ at its source. In particular all labels of $T$ are now fixed. For $r=1,2$ let $A_r,D_r$ be the products in \eqref{eq:walk-label} for $e_r$, and let $\zeta_r$ be the label of $\xi$ or $\eta$ on $e_r$. Let $\bar l_r\in\Out(S)$ be the image of $A_r\zeta_rD_r$. Since $I\lhd\Aut(S)$, the map $z\mapsto A_rzD_r$ is a bijection from $\zeta_rI$ onto the coset $\bar l_r$.

Let $E=\gen{\bar l_1,\bar l_2}\le\Out(S)$ and let $L\le\Aut(S)$ be its preimage, so that $I\le L$ and $L/I=E$. By \Cref{lem:centralizer}(b), $I$ is the unique minimal normal subgroup of $L$, and $L$ is not cyclic because it contains $I\cong S$. By \Cref{thm:lm}, $d(L)=\max\{2,d(E)\}=2$. By \Cref{thm:gaschutz} there are $l_1\in\bar l_1$ and $l_2\in\bar l_2$ with $L=\gen{l_1,l_2}$. (This includes the cases where $\bar l_1=\bar l_2$ or where one of them is trivial.) Give $e_r$ the label $z_r=A_r^{-1}l_rD_r^{-1}\in\zeta_rI$, for $r=1,2$.

Now every $x$-label has been chosen, and lies in the corresponding coset of $I$, so this defines $x\in\xi B$. Every $y$-label except those on $\mathcal R$ has been chosen in the corresponding coset. By \Cref{lem:closure} there are labels on $\mathcal R$ and an element $g\in\gamma B$ with $y=gxg^{-1}$. The walks $\omega_{e_1}$ and $\omega_{e_2}$ use only edges of $T$ and their own edge $e_r$, and none of these lies in $\mathcal R$. So by \eqref{eq:walk-label},
\[
 \omega_{e_r}(x,y)_1=A_rz_rD_r=l_r\qquad(r=1,2).
\]
The elements $\omega_{e_1}(x,y)$ and $\omega_{e_2}(x,y)$ lie in $\gen{x,y}\cap W_1$, and the homomorphism $w\mapsto w_1$ maps the subgroup they generate onto $\gen{l_1,l_2}=L\supseteq I$. Hence $|S|=|I|$ divides $|L|$, which divides $|\gen{x,y}\cap W_1|$ and therefore $|\gen{x,y}|$.
\end{proof}

\section{Proof of the theorem}\label{sec:proof}

\begin{proof}[Proof of \Cref{thm:main} and \Cref{cor:psm}]
By \Cref{lem:reduction} it suffices to prove \Cref{cor:psm}. Suppose it fails, and let $G$ be a counterexample of least order, as in \Cref{sec:radical}. Let $N_1,\dots,N_r$, $x_0,g_0,y_0$ and $\rho_i$ be as in \Cref{sec:socle}. For each $i$ we choose $a_i,h_i\in\Aut(N_i)$ with
\begin{equation}\label{eq:block-choice}
 \rho_i(x_0)^{-1}a_i,\ \rho_i(g_0)^{-1}h_i\in\Inn(N_i),\qquad |S_i|\ \text{divides}\ \bigl|\gen{a_i,\,h_ia_ih_i^{-1}}\bigr|.
\end{equation}

\emph{Blocks with $k_i=1$.} Then $N_i=S_i$ and $L_i=\rho_i(G)$ satisfies $\Inn(S_i)\le L_i\le\Aut(S_i)$. As a quotient of $G$, $L_i$ is prime-spectrum minimal by \Cref{lem:reduction}(c). The composite map $G\to L_i\to L_i/\Inn(S_i)$ is surjective and kills $M$, because $\rho_i(M)=\Inn(S_i)$. By \Cref{prop:almostsimple} the group $L_i/\Inn(S_i)$ is cyclic, so the images of $x_0$ and $y_0=g_0x_0g_0^{-1}$ in it coincide. By \eqref{eq:quotient-pair} these images generate it, so the image $c$ of $x_0$ generates $L_i/\Inn(S_i)$. Let $d$ be the image of $g_0$. \Cref{prop:almostsimple} gives $a_i\in c$ and $h_i\in d$ with $\gen{a_i,h_ia_ih_i^{-1}}=L_i\supseteq\Inn(S_i)$, so \eqref{eq:block-choice} holds.

\emph{Blocks with $k_i\ge2$.} Identify $N_i$ with $S_i^{k_i}$, so that $\Aut(N_i)=W$ and $\Inn(N_i)=B$ in the notation of \Cref{sec:wreath}. Put $\xi=\rho_i(x_0)$ and $\gamma=\rho_i(g_0)$, so that $\gamma\xi\gamma^{-1}=\rho_i(y_0)$. By \Cref{lem:transitive} the permutations $\hat\xi$ and $\widehat{\rho_i(y_0)}$ generate a transitive group. \Cref{prop:wreath} gives $a_i\in\xi B$ and $h_i\in\gamma B$ satisfying \eqref{eq:block-choice}.

\emph{Assembly.} Since $\rho_i(N_i)=\Inn(N_i)$, there are $u_i,v_i\in N_i$ with $\rho_i(u_i)=\rho_i(x_0)^{-1}a_i$ and $\rho_i(v_i)=\rho_i(g_0)^{-1}h_i$. Put
\[
 x=x_0\,u_1\cdots u_r,\qquad g=g_0\,v_1\cdots v_r,\qquad y=gxg^{-1},\qquad H=\gen{x,y}.
\]
Since $\rho_i(N_j)=1$ for $j\ne i$, we get $\rho_i(x)=\rho_i(x_0)\rho_i(u_i)=a_i$ and $\rho_i(g)=h_i$, hence $\rho_i(y)=h_ia_ih_i^{-1}$ and $\rho_i(H)=\gen{a_i,h_ia_ih_i^{-1}}$. By \eqref{eq:block-choice}, $|S_i|$ divides $|\rho_i(H)|$, which divides $|H|$. Also $x\in x_0M$ and $y\in y_0M$, so $HM=G$ by \eqref{eq:quotient-pair}, and $|G/M|=|H/(H\cap M)|$ divides $|H|$. By \Cref{lem:socle}(a),
\[
 \pi(G)=\pi(G/M)\cup\bigcup_{i=1}^r\pi(S_i)\subseteq\pi(H)\subseteq\pi(G).
\]
So $\pi(H)=\pi(G)$, and $H=G$ by prime-spectrum minimality. Then $G$ is generated by the conjugate elements $x$ and $y$, which contradicts the choice of $G$.
\end{proof}

\paragraph{Remark.}
From each block the proof needs only that $|S_i|$ divides the order of the subgroup generated by the lifts. This is what allows the blocks with $k_i\ge2$ to be handled by the label argument of \Cref{sec:wreath}, without determining the subgroup that the lifts generate. Prime-spectrum minimality then does the rest.

\section{Formal verification in Lean}\label{sec:lean}
\Cref{thm:main} has also been proved in the Lean~4 proof assistant \cite{Lean4} with the Mathlib library \cite{Mathlib}, from four of the published results of \Cref{sec:prelim}, which enter as explicit hypotheses. The development has about $2{,}000$ lines, uses version~4.24.0 of Lean and the matching release of Mathlib, and is available at \coderepo.

\emph{Statements.} The four hypotheses are the following statements, for every nonabelian finite simple group $S$ and every finite group $L$:
\begin{enumerate}
\item $\Out(S)$ is soluble (\Cref{thm:schreier});
\item a subgroup of $\Aut(S)$ whose order is prime to $|S|$ is cyclic (the case of \Cref{thm:coprime} used in \Cref{prop:almostsimple});
\item if $L$ is nonabelian and every proper quotient of $L$ is cyclic, then $L$ is generated by two conjugate elements (\Cref{thm:bgh});
\item for $\bar u,\bar v\in\Out(S)$, the preimage of $\gen{\bar u,\bar v}$ in $\Aut(S)$ is generated by two elements. This is \Cref{thm:lm} combined with \Cref{lem:centralizer}, as in the proof of \Cref{prop:wreath}.
\end{enumerate}
From these, Lean proves that every finite group $G$ has elements $x,g$ such that $\gen{x,gxg^{-1}}$ and $G$ have the same prime divisors of their orders. This is \Cref{thm:main} with $a=x$ and $b=gxg^{-1}$. The conclusion is also restated with the quantifiers written out, so that $x$ and $g$ are chosen before the prime. For soluble groups the same conclusion is proved with no hypotheses.

\emph{Proof.} The development proves Gasch\"utz's lemma (\Cref{thm:gaschutz}) by counting generating tuples, the conjugacy of complements of an abelian normal Hall subgroup used in \Cref{prop:radical}, and the part of Hall's theorem used in \Cref{prop:almostsimple}. The existence part of \Cref{thm:sz} is taken from Mathlib. The formal proof is organised slightly differently from \Cref{sec:proof}. Instead of lifting through the whole socle, it inducts on the order of the group and lifts a conjugate generating pair through a single minimal normal subgroup $N$. If $N$ is abelian it uses the argument of \Cref{prop:radical}. Otherwise it uses \Cref{prop:almostsimple} or \Cref{prop:wreath}, and prime-spectrum minimality finishes as in \Cref{sec:proof}. The simultaneous choice of lifts made in \Cref{sec:proof} is proved as a separate lemma. In \Cref{lem:closure} the formalisation takes the base point of each cycle to be the point immediately after the reserved edge. It also keeps the two directed edges of a $2$-cycle distinct, as in \Cref{lem:tree}.

\emph{Trust.} The main theorem, the soluble case and the two lifting steps of \Cref{prop:wreath,sec:proof} depend only on the three standard axioms of Lean: propositional extensionality, the axiom of choice and the soundness of quotients. The development contains no unproved steps, declares no further axioms and uses no compiled evaluation. So the formal proof relies only on the Lean kernel and on the four inputs above. A proof in Lean without hypotheses would also need formal proofs of those inputs. The first three depend on the classification of finite simple groups.

\begin{thebibliography}{99}\small
\bibitem{BGH} T.\,C.~Burness, R.\,M.~Guralnick and S.~Harper, \emph{The spread of a finite group}, Ann.\ of Math.\ (2) \textbf{193} (2021), no.~2, 619--687; arXiv:2006.01421.
\bibitem{Gaschutz} W.~Gasch\"utz, \emph{Zu einem von B.\,H. und H.~Neumann gestellten Problem}, Math.\ Nachr.\ \textbf{14} (1955), 249--252.
\bibitem{GLS3} D.~Gorenstein, R.~Lyons and R.~Solomon, \emph{The Classification of the Finite Simple Groups, Number~3}, Math.\ Surveys Monogr.\ \textbf{40.3}, Amer.\ Math.\ Soc., Providence, RI, 1998.
\bibitem{Isaacs} I.\,M.~Isaacs, \emph{Finite Group Theory}, Grad.\ Stud.\ Math.\ \textbf{92}, Amer.\ Math.\ Soc., Providence, RI, 2008.
\bibitem{Kourovka} E.\,I.~Khukhro and V.\,D.~Mazurov (eds.), \emph{Unsolved Problems in Group Theory. The Kourovka Notebook}, Sobolev Institute of Mathematics, Novosibirsk; arXiv:1401.0300 (version~50, Question~17.125 on page~123).
\bibitem{Lean4} L.~de~Moura and S.~Ullrich, \emph{The Lean~4 theorem prover and programming language}, in: Automated Deduction -- CADE~28, Lecture Notes in Comput.\ Sci.\ \textbf{12699}, Springer, 2021, 625--635.
\bibitem{LM} A.~Lucchini and F.~Menegazzo, \emph{Generators for finite groups with a unique minimal normal subgroup}, Rend.\ Sem.\ Mat.\ Univ.\ Padova \textbf{98} (1997), 173--191.
\bibitem{Maslova} N.\,V.~Maslova, \emph{On the finite prime spectrum minimal groups}, Trudy Inst.\ Mat.\ Mekh.\ UrO RAN \textbf{21} (2015), no.~3, 222--232; English transl.\ in Proc.\ Steklov Inst.\ Math.\ \textbf{295} (2016), suppl.~1, 109--119.
\bibitem{Mathlib} The mathlib Community, \emph{The Lean mathematical library}, in: Proceedings of the 9th ACM SIGPLAN International Conference on Certified Programs and Proofs (CPP 2020), ACM, 2020, 367--381.
\bibitem{MR13} N.\,V.~Maslova and D.\,O.~Revin, \emph{Generation of a finite group with Hall maximal subgroups by a pair of conjugate elements}, Trudy Inst.\ Mat.\ Mekh.\ UrO RAN \textbf{19} (2013), no.~3, 199--206; English transl.\ in Proc.\ Steklov Inst.\ Math.\ \textbf{285} (2014), suppl.~1, S139--S145.
\bibitem{SV} B.~Sp\"ath and C.~Vallejo Rodr\'iguez, \emph{Brauer characters and coprime action}, J.\ Algebra \textbf{457} (2016), 276--311.
\end{thebibliography}

\end{document}
